\documentclass[11pt,a4paper]{article}
\usepackage[UTF8]{ctex}
\usepackage{amsmath,amssymb}
\usepackage[margin=2.5cm]{geometry}
\usepackage{booktabs}

\title{Q09 解答：能量重构——双计问题的代数消除}
\author{}
\date{}

\begin{document}
\maketitle

分子：(c) H$_2$O，碎片为 O 和 H$_2$。

数值（Ha）：$E_{\text{MF}} = -74.963$，$E_O = -75.010$，$E_{\text{MF},O} = -74.980$，$E_{H_2} = -75.005$，$E_{\text{MF},H_2} = -74.983$。

\section{证明 $E_{\text{core},f}$ 抵消}

片段能量分解为 solver 无关的 core 部分和 solver 相关的电子部分：
\begin{align}
  E_f &= E_{\text{core},f} + E_{\text{elec},f} \\
  E_{\text{MF},f} &= E_{\text{core},f} + E_{\text{HF-elec},f}
\end{align}

\begin{align}
  E_f - E_{\text{MF},f}
  &= \big(E_{\text{core},f} + E_{\text{elec},f}\big) - \big(E_{\text{core},f} + E_{\text{HF-elec},f}\big) \\
  &= \big(E_{\text{core},f} - E_{\text{core},f}\big) + \big(E_{\text{elec},f} - E_{\text{HF-elec},f}\big) \\
  &= E_{\text{elec},f} - E_{\text{HF-elec},f}
\end{align}

$E_{\text{core},f}$ 完全抵消，只剩下电子关联修正：

\begin{equation}
  \boxed{\;\Delta E_{\text{corr},f} = E_f - E_{\text{MF},f} = E_{\text{elec},f} - E_{\text{HF-elec},f}\;}
\end{equation}

\section{为什么 $\Delta E_{\text{corr},f}$ 局域且无双计}

\subsection{局域性}

$\Delta E_{\text{corr},f}$ 只涉及碎片 $f$ 的嵌入空间（碎片轨道 + 浴轨道）内的电子关联。电子关联主要是局域的——电子主要和近邻电子纠缠。跨碎片关联（色散、电荷转移）在标准 DMET 中被丢失，但\textbf{不是双计问题}——它是系统性遗漏，属于 $\varepsilon_{\text{cross}}$ 误差源。

\subsection{无双计的原因}

朴素求和 $\sum_f E_f$ 会偏大，因为相邻碎片的重叠浴轨道的 HF 贡献被重复计算。标准重构通过减去 $E_{\text{MF},f}$ 消除了这个重复：

\begin{itemize}
  \item $\sum_f E_{\text{MF},f}$ 包含了重叠浴的\textbf{平均场版本}贡献
  \item $\sum_f E_f$ 包含了重叠浴的\textbf{精确版本}贡献
  \item 取差 $\sum_f (E_f - E_{\text{MF},f})$：重叠部分的平均场被精确扣除，只替换为关联修正
\end{itemize}

\textbf{核能 $E_{NN}$ 只出现在 $E_{\text{MF}}$ 中一次}——核位置在 DMET 碎片化中固定，$E_f$ 和 $E_{\text{MF},f}$ 都只是电子能量，差值不含核能，无需按碎片分解。


\section{验证：全 HF 时 $E_{\text{total}} = E_{\text{MF}}$}

若所有碎片都用 HF 求解器，则 $E_f = E_{\text{MF},f}$（同一空间同一方法），每个差值为零：

\begin{equation}
  E_{\text{total}} = E_{\text{MF}} + \sum_f \underbrace{(E_f - E_{\text{MF},f})}_{=\,0} = E_{\text{MF}} \quad \checkmark
\end{equation}

\section{计算重构后的总能量}

\begin{equation}
  E_{\text{total}} = E_{\text{MF}} + \sum_f \big(E_f - E_{\text{MF},f}\big) = E_{\text{MF}} + \big(E_O - E_{\text{MF},O}\big) + \big(E_{H_2} - E_{\text{MF},H_2}\big)
\end{equation}

\begin{align}
  \Delta E_{\text{corr},O} &= E_O - E_{\text{MF},O} = (-75.010) - (-74.980) = -0.030 \text{ Ha} \\
  \Delta E_{\text{corr},H_2} &= E_{H_2} - E_{\text{MF},H_2} = (-75.005) - (-74.983) = -0.022 \text{ Ha}
\end{align}

\begin{align}
  E_{\text{total}} &= -74.963 + (-0.030) + (-0.022) \\
  &= -74.963 - 0.052 \\
  &= \boxed{-75.015 \text{ Ha}}
\end{align}

\section*{分块累积量重构}

\subsection*{标准重构的局限}

标准公式 $E_{\text{total}} = E_{\text{MF}} + \sum_f \Delta E_{\text{corr},f}$ 把能量当作一个整体处理，假设关联能可按碎片简单相加。但跨碎片关联（色散、电荷转移）被丢失。

\subsection*{累积量分解的基本思想}

把总能量按\textbf{连接阶数}分解为累积量：
$$E = E^{(1)} + E^{(2)} + E^{(3)} + \cdots$$

\begin{itemize}
  \item $E^{(1)}$：1-body 累积量 $\approx$ HF 贡献（单电子 + 平均场排斥）
  \item $E^{(2)}$：2-body 累积量 $\approx$ 电子对关联（关联能的主要部分）
  \item $E^{(3)}$ 及更高：三体及以上关联，通常较小可忽略
\end{itemize}

对\textbf{每一阶分别重构}：
$$E_{\text{total}} = E_{\text{MF}} + \sum_f \Big[\big(E^{(2)}_f - E^{(2)}_{\text{MF},f}\big) + \big(E^{(3)}_f - E^{(3)}_{\text{MF},f}\big) + \cdots\Big]$$

\subsection*{何时明显更精确}

\begin{itemize}
  \item \textbf{跨碎片 2-body 关联显著时}：色散力、$\pi$ 共轭体系。标准重构丢失跨碎片关联，但 2-body 累积量可通过浴轨道重叠部分捕获。
  \item \textbf{强关联体系}：N$_2$ 拉伸到 2.0 \AA{} 时，静态关联占主导，2-body 累积量的非局域贡献变大。
  \item \textbf{大基组}：虚拟轨道多，跨碎片虚拟-虚拟关联更丰富，累积量分解能更好分离局域和非局域贡献。
\end{itemize}

\textbf{代价}：需要计算每个碎片的 2-body 累积量（2-RDM 的连接部分），计算量比标准重构（只需能量）大得多。

\end{document}
